\section*{Capítulo \ref{ch:b1:series} --- Series numéricas}

\begin{solution}{exo:b1:series:1}
$\dfrac{1}{n(n+1)} = \dfrac1n - \dfrac{1}{n+1}$: telescópica,
$S_N = 1 - \frac{1}{N+1} \to 1$. Convergente, con suma $1$.

$\ln\bigl(1 - \frac{1}{n^2}\bigr) = \ln\frac{(n-1)(n+1)}{n^2} =
\ln\frac{n-1}{n} - \ln\frac{n}{n+1}$: otra vez telescópica,
$S_N = \ln\frac12 - \ln\frac{N}{N+1} \to -\ln 2$. Convergente, con
suma $-\ln 2$.

$\dfrac{3^n + 4^n}{5^n} = \bigl(\frac35\bigr)^n +
\bigl(\frac45\bigr)^n$: dos \omterm{ex:b1:series:geometric}{series geométricas} convergentes, con suma
$\frac{1}{1 - 3/5} + \frac{1}{1 - 4/5} = \frac52 + 5 = \frac{15}{2}$.
\end{solution}

\begin{solution}{exo:b1:series:2}
\omterm{thm:b1:series:ratio}{Criterio del cociente} en todos los casos.

$\frac{u_{n+1}}{u_n} = \frac{(n+1)^2}{2n^2} \to \frac12 < 1$:
convergente.

$\frac{u_{n+1}}{u_n} = \frac{(n+1)!\,n^n}{n!\,(n+1)^{n+1}} =
\bigl(\frac{n}{n+1}\bigr)^n = \bigl(1 + \frac1n\bigr)^{-n} \to
\frac1\eu < 1$: convergente.

Con el factor $2^n$: cociente $\to \frac2\eu < 1$: convergente.

Con $3^n$: cociente $\to \frac3\eu > 1$: divergente (los términos
tienden a $+\infty$).
\end{solution}

\begin{solution}{exo:b1:series:3}
Todos los términos son no negativos; úsense equivalentes
(\cref{thm:b1:series:positive}).

$\sin\frac{1}{n^2} \sim \frac{1}{n^2}$: convergente (Riemann con
$\alpha = 2$).

$1 - \cos\frac1n \sim \frac{1}{2n^2}$: convergente.

$\frac{1}{\sqrt{n(n+1)}} \sim \frac1n$: divergente.

$\frac{\ln n}{n^2} = \frac{1}{n^{3/2}}\cdot\frac{\ln n}{n^{1/2}}$ y
$\frac{\ln n}{\sqrt n} \to 0$
(\cref{prop:b1:functions:powerrules}): luego
$\frac{\ln n}{n^2} \leq \frac{1}{n^{3/2}}$ para $n$ grande:
convergente.
\end{solution}

\begin{solution}{exo:b1:series:4}
Para $n \geq 2$: $\frac{1}{n^2} \leq \frac{1}{n(n-1)} =
\frac{1}{n-1} - \frac1n$. Por tanto,
\[
\sum_{n=1}^{N} \frac{1}{n^2} \leq 1 + \sum_{n=2}^{N}
\Bigl(\frac{1}{n-1} - \frac1n\Bigr) = 1 + 1 - \frac1N < 2 :
\]
sumas parciales crecientes y acotadas por $2$: convergencia
(\cref{thm:b1:series:positive}), con suma $\leq 2$. (El valor exacto
$\frac{\pi^2}{6}$ es una celebración de segundo año.)
\end{solution}

\begin{solution}{exo:b1:series:5}
$\sum \frac{(-1)^n}{\sqrt n}$: alternada con
$\frac{1}{\sqrt n} \downarrow 0$: convergente
(\cref{thm:b1:series:alternating}); no absolutamente
($\alpha = \frac12 \leq 1$).

$\sum \frac{(-1)^n}{n + (-1)^n}$: la sucesión
$\frac{1}{n + (-1)^n}$ \emph{no} es decreciente ($\frac{1}{n+1}$ y
después $\frac{1}{n}$ se alternan mal), de modo que el criterio no se
aplica directamente. Desarróllese:
\[
\frac{(-1)^n}{n + (-1)^n}
= \frac{(-1)^n}{n}\cdot\frac{1}{1 + \frac{(-1)^n}{n}}
= \frac{(-1)^n}{n} - \frac{1}{n^2} +
O\Bigl(\frac{1}{n^3}\Bigr):
\]
la primera \omterm{def:b1:series:def}{serie} converge (alternada), $\sum \frac{1}{n^2}$
converge y el $O(n^{-3})$ converge absolutamente: la suma de tres
\omterm{def:b1:series:def}{series} convergentes converge.

$\sin\bigl(\pi\sqrt{n^2+1}\bigr)$: escríbase
$\sqrt{n^2 + 1} = n + \frac{1}{2n} + \varepsilon_n$ con
$\varepsilon_n = O(n^{-3})$; entonces, por la $\pi$-periodicidad de
$\sin$ salvo el signo,
\[
\sin\bigl(\pi\sqrt{n^2+1}\bigr)
= (-1)^n \sin\Bigl(\frac{\pi}{2n} + \pi\varepsilon_n\Bigr) .
\]
Póngase $\theta_n = \frac{\pi}{2n} + \pi\varepsilon_n$ y
$a_n = \sin\theta_n$. Para $n$ grande,
$\theta_n \in \intoo{0}{\frac\pi2}$ y
\[
\theta_n - \theta_{n+1} = \frac{\pi}{2n(n+1)} +
\pi(\varepsilon_n - \varepsilon_{n+1})
= \frac{\pi}{2n^2} + O\Bigl(\frac{1}{n^3}\Bigr) > 0
\]
a la larga, de modo que $(\theta_n)$ decrece hacia $0$; y como $\sin$
es creciente en $\intcc{0}{\frac\pi2}$, también $(a_n)$ decrece hacia
$0$. Se aplica el criterio de las alternadas: convergente --- no
absolutamente, ya que $a_n \sim \frac{\pi}{2n}$.
\end{solution}

\begin{solution}{exo:b1:series:6}
$f(t) = \frac{1}{t(\ln t)^\alpha}$ es positiva, \omterm{def:b1:continuity:continuous}{continua} y
decreciente en $\intco{2}{+\infty}$. Sustituyendo $u = \ln t$:
\[
\int_2^x \frac{\dd t}{t(\ln t)^\alpha}
= \int_{\ln 2}^{\ln x} \frac{\dd u}{u^\alpha},
\]
acotada cuando $x \to \infty$ si y solo si $\alpha > 1$ (el cálculo
del~\cref{thm:b1:series:riemann}). Por comparación con \omterm{thm:b1:integration:def}{integrales}:
convergencia si y solo si $\alpha > 1$. (Estas \omterm{def:b1:series:def}{series} \emph{de tipo
Bertrand} muestran lo fina que es la \omterm{def:b1:topology:closure}{frontera} de la convergencia:
$n\ln n$ diverge y $n(\ln n)^{1.01}$ converge.)
\end{solution}

\begin{solution}{exo:b1:series:7}
Por las cotas de tangencia del~\cref{exo:b1:derivative:3}, reescritas
mediante desarrollos: para $x = \frac1n \in \intoc{0}{1}$,
Taylor--Lagrange para $\ln(1+x)$ al orden $1$ da
$\ln(1 + x) = x - \frac{x^2}{2(1 + c)^2}$ para algún
$c \in \intoo{0}{x}$, luego
\[
0 \leq u_n = \frac1n - \ln\Bigl(1 + \frac1n\Bigr) \leq
\frac{1}{2n^2} .
\]
Comparación con la \omterm{def:b1:series:def}{serie} de Riemann: $\sum u_n$ converge. Y su suma
parcial telescopa los logaritmos:
\[
\sum_{n=1}^{N} u_n = H_N - \ln(N+1)
\]
(puesto que $\sum_{n\leq N} \ln\frac{n+1}{n} = \ln(N+1)$). Así,
$H_N - \ln(N+1)$ converge; y añadiendo $\ln\frac{N+1}{N} \to 0$, la
sucesión $H_N - \ln N$ converge. Su límite es $\gamma \approx 0.5772$.
\end{solution}

\begin{solution}{exo:b1:series:8}
$\dfrac{1}{n(n+2)} = \dfrac{1/2}{n} - \dfrac{1/2}{n+2}$: la suma
parcial telescopa con un retardo de $2$,
\[
S_N = \frac12\Bigl(1 + \frac12 - \frac{1}{N+1} -
\frac{1}{N+2}\Bigr) \longrightarrow \frac34 .
\]

$\sum \frac{n}{2^n}$: para $\abs x < 1$, derivando la suma geométrica
finita y pasando al límite (todas las \omterm{def:b1:series:def}{series} de aquí convergen
absolutamente, \omterm{thm:b1:series:ratio}{criterio del cociente}): de
$\sum_{n\geq0} x^n = \frac{1}{1-x}$ se obtiene, por cálculo directo
con sumas parciales,
\[
\sum_{n=1}^{N} n x^{n-1}
= \frac{1 - (N+1)x^N + N x^{N+1}}{(1 - x)^2}
\xrightarrow[N\to\infty]{} \frac{1}{(1-x)^2}
\quad (\abs x < 1),
\]
(los términos de \omterm{def:b1:topology:closure}{frontera} cumplen $N x^N \to 0$). En $x = \frac12$:
$\sum_{n\geq1} n\bigl(\frac12\bigr)^{n-1} = 4$, luego
$\sum_{n\geq0} \frac{n}{2^n} = \frac12 \times 4 = 2$.
\end{solution}

\begin{solution}{exo:b1:series:9}
Agrúpense los términos de $\sum u_n$ en paquetes entre potencias de
$2$. Paquetes superiores: para $2^k \leq n < 2^{k+1}$ hay $2^k$
términos, cada uno $\leq u_{2^k}$:
\[
\sum_{n=1}^{2^{K+1}-1} u_n
= \sum_{k=0}^{K} \sum_{n=2^k}^{2^{k+1}-1} u_n
\leq \sum_{k=0}^{K} 2^k u_{2^k} .
\]
Paquetes inferiores: cada término del mismo paquete es
$\geq u_{2^{k+1}}$, luego
$\sum_{n=2^k}^{2^{k+1}-1} u_n \geq 2^k u_{2^{k+1}} = \frac12
\cdot 2^{k+1} u_{2^{k+1}}$, de donde
\[
\sum_{n=1}^{2^{K+1}-1} u_n \geq \frac12 \sum_{k=1}^{K+1} 2^{k}
u_{2^{k}} .
\]
Las dos comparaciones de sumas parciales funcionan en los dos
sentidos (términos no negativos, \cref{thm:b1:series:positive}): las
dos \omterm{def:b1:series:def}{series} tienen la misma naturaleza.

Riemann: $u_n = n^{-\alpha}$ da $2^k u_{2^k} = 2^{k(1-\alpha)}$, una
\omterm{ex:b1:series:geometric}{serie geométrica}, convergente si y solo si $2^{1 - \alpha} < 1$, es
decir, si y solo si $\alpha > 1$. Bertrand
(\cref{exo:b1:series:6}): $u_n = \frac{1}{n(\ln n)^\alpha}$ da
$2^k u_{2^k} = \frac{1}{(k\ln 2)^\alpha}$, una \omterm{def:b1:series:def}{serie} de Riemann en
$k$: convergente si y solo si $\alpha > 1$.
\end{solution}

\begin{solution}{exo:b1:series:10}
Identidad geométrica finita de razón $-t^2$:
\[
\frac{1}{1 + t^2} = \sum_{k=0}^{n-1} (-1)^k t^{2k}
+ \frac{(-1)^n t^{2n}}{1 + t^2} .
\]
Intégrese en $\intcc{0}{1}$ (el miembro izquierdo se integra a
$\arctan 1 = \frac\pi4$, \cref{prop:b1:functions:arcderiv}):
\[
\frac{\pi}{4} = \sum_{k=0}^{n-1} \frac{(-1)^k}{2k+1}
+ (-1)^n \int_0^1 \frac{t^{2n}}{1+t^2}\,\dd t,
\qquad
0 \leq \int_0^1 \frac{t^{2n}}{1+t^2}\,\dd t \leq \int_0^1 t^{2n}\dd
t = \frac{1}{2n+1} .
\]
Haciendo $n \to \infty$ se demuestra la fórmula, y la cota anterior
sobre la \omterm{thm:b1:integration:def}{integral} es exactamente la cota del resto: tras sumar hasta
$N$ (es decir, $n = N + 1$ términos), $\abs{R_N} \leq \frac{1}{2N + 3}$.
\end{solution}

\begin{solution}{exo:b1:series:11}
$n^{1/n} = \eu^{\frac{\ln n}{n}} \to \eu^0 = 1$
(\cref{prop:b1:functions:powerrules}). Por tanto,
\[
\frac{1}{n^{1 + 1/n}} = \frac{1}{n}\,\eu^{-\frac{\ln n}{n}}
\sim \frac{1}{n} ,
\]
y el criterio de los equivalentes (\cref{thm:b1:series:positive})
compara con la \omterm{def:b1:series:def}{serie} armónica divergente: \emph{divergente}, aunque
todo exponente $1 + \frac1n$ supere a $1$. El criterio de Riemann
concierne a un exponente \emph{fijo} $\alpha$; un exponente que se
desliza hacia $1$ puede perder todo su margen, como aquí.
\end{solution}

\begin{solution}{exo:b1:series:12}
Sea $\varepsilon > 0$. Por el criterio de Cauchy para la \omterm{def:b1:series:def}{serie}
convergente (\cref{thm:b1:seq:complete} aplicado a las sumas
parciales), hay $N$ con $\sum_{k=n+1}^{2n} u_k \leq \varepsilon$ para
$n \geq N$. Por monotonía, cada uno de esos $n$ términos es
$\geq u_{2n}$:
\[
n\,u_{2n} \leq \sum_{k=n+1}^{2n} u_k \leq \varepsilon
\quad\Longrightarrow\quad 2n\,u_{2n} \leq 2\varepsilon ,
\]
y para los índices impares,
$(2n+1)\,u_{2n+1} \leq (2n+1)\,u_{2n} \leq
2\bigl(2n\,u_{2n}\bigr) \leq 4\varepsilon$ para $n \geq N$: en las
dos paridades, $n u_n \to 0$.

Recíproco falso: $u_n = \frac{1}{n\ln n}$ tiene
$n u_n = \frac{1}{\ln n} \to 0$ y, sin embargo, la \omterm{def:b1:series:def}{serie} diverge
(\cref{exo:b1:series:6}, $\alpha = 1$). Monotonía necesaria: sea
$u_n = \frac1n$ cuando $n$ es un cuadrado perfecto y $u_n = 2^{-n}$
en los demás casos: la \omterm{def:b1:series:def}{serie} converge (los términos cuadrados suman
como $\sum \frac{1}{k^2}$, y el resto, geométricamente), pero
$n u_n = 1$ a lo largo de los cuadrados.
\end{solution}

\begin{solution}{pb:b1:series:1}
\textbf{1.} $a_{n+1} - a_n = \frac{1}{n+1} - \ln\frac{n+1}{n}
\leq 0$ porque $\ln(1 + \frac1n) \geq \frac{1/n}{1 + 1/n} =
\frac{1}{n+1}$; y $b_{n+1} - b_n = \frac{1}{n+1} -
\ln\frac{n+2}{n+1} \geq 0$ porque
$\ln(1 + \frac{1}{n+1}) \leq \frac{1}{n+1}$. Su hueco
$a_n - b_n = \ln(1 + \frac1n) \to 0$: adyacentes
(\cref{thm:b1:seq:adjacent}), con límite común
$\lim a_n = \gamma$ (\cref{exo:b1:series:7}), y
$b_n \leq \gamma \leq a_n$.

\textbf{2.} $b_{10} = 2.928968 - \ln 11 = 0.5311$ y
$a_{10} = 2.928968 - \ln 10 = 0.6264$: luego
$\gamma \in \intcc{0.5311}{0.6264}$. El hueco vale
$\ln 1.1 \approx 0.095$ y se encoge como $\frac1n$: cuatro decimales
($\text{hueco} \leq 10^{-4}$) exigirían $n \approx 10^4$ --- los
encuadres son correctos, pero lentos.

\textbf{3.} Telescopando $a_n - a_{m+1} = \sum_{k=n}^{m} w_k$ y
haciendo $m \to \infty$: $a_n - \gamma = \sum_{k\geq n} w_k$ (límite
de sumas parciales). Además,
\[
w_k = \int_k^{k+1} \frac{\dd t}{t} - \frac{1}{k+1}
= \int_k^{k+1} \Bigl(\frac1t - \frac{1}{k+1}\Bigr)\dd t
= \int_k^{k+1} \frac{k + 1 - t}{t\,(k+1)}\,\dd t .
\]
En $\intcc{k}{k+1}$: $\frac{1}{(k+1)^2} \leq \frac{1}{t(k+1)}
\leq \frac{1}{k(k+1)}$, y $\int_k^{k+1}(k + 1 - t)\dd t =
\frac12$: por tanto,
$\frac{1}{2(k+1)^2} \leq w_k \leq \frac{1}{2k(k+1)}$.

\textbf{4.} Por arriba:
$\sum_{k \geq n} \frac{1}{2k(k+1)} =
\frac12\sum_{k\geq n}\bigl(\frac1k - \frac{1}{k+1}\bigr) =
\frac{1}{2n}$ (telescopaje). Por abajo:
$\frac{1}{2(k+1)^2} \geq \frac{1}{2(k+1)(k+2)}$, cuya suma telescopa
hasta $\frac{1}{2(n+1)}$. Con la pregunta 3:
\[
\frac{1}{2(n+1)} \leq H_n - \ln n - \gamma \leq \frac{1}{2n} .
\]

\textbf{5.} Réstese $\frac{1}{2n}$:
$\gamma_n - \gamma \in \intcc{\frac{1}{2(n+1)} - \frac{1}{2n}}{0} =
\intcc{-\frac{1}{2n(n+1)}}{0}$: la estimación corregida es exacta
salvo $O\bigl(\frac{1}{n^2}\bigr)$, y siempre por debajo.

\textbf{6.} $\gamma_{100} = 5.1873775 - \ln 100 - 0.005 =
0.5772073$, con
$0 \leq \gamma - \gamma_{100} \leq \frac{1}{20200} < 5\cdot10^{-5}$:
por tanto, $0.577207 \leq \gamma \leq 0.577257$, es decir,
$\gamma = 0.5772 \pm 5\cdot10^{-5}$ (valor verdadero
$0.5772156\dots$) --- cuatro decimales certificados con cien
términos, frente a los diez mil de la pregunta 2.

\textbf{7.} Primer dividendo:
\[
H_{2m} - H_m = \Bigl(\ln 2m + \gamma + \frac{1}{4m}\Bigr)
- \Bigl(\ln m + \gamma + \frac{1}{2m}\Bigr) +
O\Bigl(\frac{1}{m^2}\Bigr)
= \ln 2 - \frac{1}{4m} + O\Bigl(\frac{1}{m^2}\Bigr).
\]
Segundo: los inversos pares hasta $2m$ suman $\frac12 H_m$, luego
$\sum_{j=1}^{m}\frac{1}{2j-1} = H_{2m} - \frac12 H_m = \frac12
\ln m + \ln 2 + \frac\gamma2 + o(1)$, y
$\sum_{j=1}^m \frac{1}{2j} = \frac12\ln m + \frac\gamma2 + o(1)$.

\textbf{8.} Separando dos veces los términos pares:
$\sum_{k=1}^{2m}\frac{(-1)^{k-1}}{k} = H_{2m} - 2\cdot\frac12
H_m = H_{2m} - H_m$. Por la pregunta 7, esto vale
$\ln 2 - \frac{1}{4m} + O(m^{-2})$: la suma es $\ln 2$
(otra vez el~\cref{ex:b1:series:ln2}) \emph{con} su velocidad.

\textbf{9.} Para $N = 2m$:
$S - S_N = \frac{1}{4m} + O(m^{-2}) = \frac{1}{2N} + O(N^{-2})$.
Para $N = 2m + 1$: $S_{N} = S_{2m} + \frac{1}{2m+1}$, luego
\[
S - S_N = \Bigl(\frac{1}{4m} - \frac{1}{2m+1}\Bigr) +
O\Bigl(\frac{1}{m^2}\Bigr)
= -\frac{1}{4m} + O\Bigl(\frac{1}{m^2}\Bigr)
= -\frac{1}{2N} + O\Bigl(\frac{1}{N^2}\Bigr).
\]
En los dos casos, $S - S_N \sim \frac{(-1)^N}{2N}$: la mitad de la
cota del peor caso $a_{N+1}$, y con signo conocido y alternante.

\textbf{10.} Promediar mata el término dominante oscilante:
\[
S - \tilde S_N = \frac{(S - S_N) + (S - S_{N+1})}{2}
= \frac{(-1)^N}{2}\Bigl(\frac{1}{2N} - \frac{1}{2(N+1)}\Bigr)
+ O\Bigl(\frac{1}{N^2}\Bigr) = O\Bigl(\frac{1}{N^2}\Bigr).
\]
Numéricamente: $S_{10} = 0.645635$, $S_{11} = 0.736544$,
$\tilde S_{10} = 0.691089$ y $\ln 2 = 0.693147$: el error cae de
$4.8\cdot10^{-2}$ a $2.1\cdot10^{-3}$ --- una suma, veinte veces
mejor.

\textbf{11.} El error alternado cambia de signo en cada paso, de modo
que dos sumas parciales consecutivas cabalgan sobre el límite y su
media cancela el término de primer orden; el error del encuadre
$H_n - \ln n - \gamma \approx \frac{1}{2n}$ tiene signo constante, y
ningún promedio sobre $n$ puede cancelarlo. Promediar los dos
\emph{encuadres} \emph{sí} ayuda:
$\frac{a_n + b_n}{2} = H_n - \ln\sqrt{n(n+1)}$ y, como
$\ln\sqrt{n(n+1)} = \ln\bigl(n + \frac12\bigr) + O(n^{-2})$,
\[
H_n - \ln\Bigl(n + \frac12\Bigr)
= \Bigl(H_n - \ln n - \frac{1}{2n}\Bigr) + \frac{1}{8n^2} +
O\Bigl(\frac{1}{n^3}\Bigr) = \gamma + O\Bigl(\frac{1}{n^2}\Bigr)
\]
(pregunta 5 y
$\ln(1 + \frac{1}{2n}) = \frac{1}{2n} - \frac{1}{8n^2} + O(n^{-3})$).
Comprobación en $n = 10$: $H_{10} - \ln 10.5 = 0.57759$, ya a menos
de $4\cdot10^{-4}$ de $\gamma$.

\textbf{12.} $\sum_{j\leq m} \frac{1}{2j-1} \geq \sum_{j \leq m}
\frac{1}{2j} = \frac12 H_m \to +\infty$: tanto la parte positiva
como la negativa de la \omterm{def:b1:series:def}{serie} armónica alternada divergen.

\textbf{13.} Escríbase $u_n^\pm = \frac{\abs{u_n} \pm u_n}{2} \geq
0$, de modo que $u_n = u_n^+ - u_n^-$ y
$\abs{u_n} = u_n^+ + u_n^-$. Si $\sum u_n^+$ convergiera, entonces
$\sum u_n^- = \sum (u_n^+ - u_n)$ convergería (diferencia de \omterm{def:b1:series:def}{series}
convergentes), y por tanto también $\sum \abs{u_n}$: contradicción
con la convergencia condicional. Por simetría, las dos $\sum u_n^\pm$
divergen (a $+\infty$): un reservorio infinito de masa positiva y de
masa negativa.

\textbf{14.} Sean $S = \sum u_n$, $\varepsilon > 0$ y $N$ con
$\sum_{n > N} \abs{u_n} \leq \varepsilon$ (criterio de Cauchy para
$\sum\abs{u_n}$). Sea $M_0$ lo bastante grande para que
$\sigma(\{0, \dots, M_0\}) \supseteq \{0, \dots, N\}$. Para
$M \geq M_0$, la diferencia
$\sum_{m \leq M} u_{\sigma(m)} - \sum_{n \leq N} u_n$ es una suma
finita de términos $u_n$ \emph{distintos} con $n > N$, luego de valor
absoluto $\leq \varepsilon$; y también
$\abs{S - \sum_{n\leq N} u_n} \leq \varepsilon$. Así, las sumas
parciales reordenadas acaban estando a menos de $2\varepsilon$ de
$S$: $\sum u_{\sigma(n)} = S$. La \omterm{thm:b1:series:absolute}{convergencia absoluta} es a prueba
de reordenaciones.

\textbf{15.} Cada fase del procedimiento voraz termina tras finitos
términos, porque los términos positivos (respectivamente, negativos)
que quedan tienen ya por sí solos sumas parciales divergentes
(pregunta 12): la suma en curso tiene que acabar cruzando $t$. El
procedimiento alterna, por tanto, infinitas fases finitas,
consumiendo los términos positivos en orden y los negativos en orden:
cada término se usa exactamente una vez --- una reordenación. Tras el
primer cruce, entre dos cruces consecutivos las sumas parciales se
mueven monótonamente hacia $t$, y en un cruce se pasan a lo sumo en
el término recién añadido; y como los términos usados en el
$j$-ésimo cruce tienen índice al menos $j$ dentro de su clase, esos
excesos tienden a $0$. Luego las sumas parciales convergen a $t$:
todo número real es la suma de alguna reordenación.

\textbf{16.} Las casillas positivas reciben $\frac{1}{2j-1}$ para
$j = 1, 2, \dots$ en orden, y las negativas, $\frac{1}{2j}$ en
orden: cada término de la \omterm{def:b1:series:def}{serie} armónica alternada aparece
exactamente una vez. Para $(p, q) = (1, 2)$, los bloques son
$\bigl(1, -\frac12, -\frac14\bigr)$, $\bigl(\frac13, -\frac16,
-\frac18\bigr)$, $\bigl(\frac15, -\frac1{10},
-\frac1{12}\bigr)$, \dots\ --- la \omterm{def:b1:series:def}{serie} escrita antes.

\textbf{17.} Como $\frac{1}{4k-2} = \frac12\cdot\frac{1}{2k-1}$:
\[
\frac{1}{2k-1} - \frac{1}{4k-2} - \frac{1}{4k}
= \frac12\,\frac{1}{2k-1} - \frac12\,\frac{1}{2k}
= \frac12\Bigl(\frac{1}{2k-1} - \frac{1}{2k}\Bigr).
\]
Sumando sobre $k = 1, \dots, K$:
$T_{3K} = \frac12 \sum_{k=1}^{K}\bigl(\frac{1}{2k-1} -
\frac{1}{2k}\bigr) = \frac12 S_{2K}$: en cada tercera suma parcial,
la \omterm{def:b1:series:def}{serie} reordenada es \emph{exactamente} la mitad de la original.

\textbf{18.} Tras $K$ bloques completos, la suma parcial reordenada
es $\sum_{j=1}^{pK}\frac{1}{2j-1} - \sum_{j=1}^{qK}\frac{1}{2j}$, y
la pregunta 7 la evalúa:
\[
\Bigl(\frac{\ln(pK)}{2} + \ln 2 + \frac\gamma2\Bigr)
- \Bigl(\frac{\ln(qK)}{2} + \frac\gamma2\Bigr) + o(1)
= \ln 2 + \frac12\ln\frac pq + o(1) :
\]
las $\gamma$ se cancelan, los $\ln K$ se cancelan y sobrevive la
razón $\frac pq$.

\textbf{19.} Una suma parcial dentro del bloque $K + 1$ difiere de la
suma tras $K$ bloques en a lo sumo $p + q$ términos, cada uno de
valor absoluto aproximadamente $\leq \frac{1}{2qK}$, es decir, en
$O\bigl(\frac1K\bigr) \to 0$: toda la sucesión de sumas parciales
tiene el mismo límite $\ln 2 + \frac12\ln\frac pq$. Para $(1, 2)$:
$\ln 2 + \frac12\ln\frac12 = \frac{\ln 2}{2} = 0.34657\dots$, y en
efecto $T_9 = 0.30833$ avanza hacia él: por la pregunta 17,
$T_{3K} = \frac12 S_{2K}$ converge con exactamente la mitad del error
de la armónica alternada. Los mismos términos, la mitad de la suma.

\textbf{20.} $(1,1)$: $\ln 2 + \frac12\ln 1 = \ln 2$ --- el orden
original, coherencia. $(2,1)$: $\frac32\ln 2 \approx 1.0397$. El menú
$(p,q)$ alcanza exactamente la familia numerable y \omterm{def:b1:topology:dense}{densa}
$\ln 2 + \frac12\ln r$, $r \in \Q_{>0}$; la receta voraz de Riemann
(pregunta 15) alcanza \emph{todos} los reales. La estructura compra
fórmulas; la voracidad compra totalidad.

\textbf{21.} Las sumas parciales telescopan:
$\sum_{k=1}^{N} \bigl(\frac1k - \ln\frac{k+1}{k}\bigr) =
H_N - \ln(N+1) = b_N \to \gamma$: la \omterm{def:b1:series:def}{serie} del~\cref{exo:b1:series:7}
suma exactamente la \omterm{pb:b1:series:1}{constante de Euler}.

\textbf{22.} Voraz para $t = 1$: el primer término positivo lleva la
suma exactamente a $1$, no más allá, de modo que se toma un segundo
positivo para cruzar: $1, \frac13$ (suma $1.3333 > 1$), después
$-\frac12$ ($0.8333$), $\frac15$ ($1.0333$), $-\frac14$
($0.7833$), $\frac17, \frac19$ ($1.0373$), $-\frac16$
($0.8706$), $\frac1{11}, \frac1{13}$ ($1.0384$), $-\frac18$
($0.9134$), $\frac1{15}$ ($0.9801$), \dots\ --- las sumas respiran en
torno a $1$ con amplitud cada vez menor, y ahora hacen falta dos
positivos por ciclo, porque los negativos son mayores.

\textbf{23.} Constrúyanse bloques: en la etapa $j$, añádanse
suficientes términos positivos sin usar para subir la suma parcial al
menos en $1$ (es posible: los términos positivos restantes tienen
sumas divergentes, pregunta 12) y añádase después el único término
negativo $-\frac{1}{2j}$. Todo término positivo acaba usándose (cada
etapa usa al menos uno) y todo negativo también (uno por etapa): una
reordenación. Cada etapa cambia la suma en $\geq 1 - \frac{1}{2j}
\geq \frac12$: las sumas parciales superan $\frac{j}{2}$ tras la
etapa $j$, y los incrementos dentro de una etapa son positivos salvo
el último, acotado por $\frac{1}{2j} \to 0$: divergencia a $+\infty$.

\textbf{24.} Por la ley,
$H_N \geq 20 \iff \ln N \geq 20 - \gamma - \frac{1}{2N} +
O(N^{-2})$: el umbral $N^*$ cumple $\ln N^* = 20 - \gamma + o(1)$, es
decir, $N^* = \eu^{20 - \gamma}(1 + o(1)) \approx \eu^{19.4228}
\approx 2.72\cdot10^{8}$ --- dentro de la ventana tosca
$\intcc{1.8\cdot10^8}{4.9\cdot10^8}$
del~\cref{ex:b1:series:harmonicstack}, y fijado por $\gamma$.

\textbf{25.} (i) En $H_n = \ln n + \gamma + \frac{1}{2n} +
O(n^{-2})$: el $\ln n$ es la \omterm{thm:b1:integration:def}{integral}, $\gamma$ es el precio de
sustituir una suma por una \omterm{thm:b1:integration:def}{integral} (una constante del análisis
genuinamente nueva) y $\frac{1}{2n}$ es la primera corrección --- la
sombra del trapecio. (ii) La convergencia condicional se apoya en la
cancelación entre dos reservorios infinitos (pregunta 13), de modo
que reordenar reajusta los pesos de los reservorios; la
\omterm{thm:b1:series:absolute}{convergencia absoluta} tiene masa total finita, y la estimación de
la cola de la pregunta 14 es ciega al orden. (iii) Las sumas $(p,q)$
se \emph{calcularon}: la ley de $\gamma$ convirtió cada suma parcial
reordenada en $\ln 2 + \frac12\ln\frac pq + o(1)$, cancelándose el
propio $\gamma$ --- un ejercicio de contabilidad asintótica, no un
argumento abstracto. (iv) A continuación: los productos y la
sumabilidad incondicional para las \omterm{def:b1:series:def}{series} de potencias en el volumen
del segundo año; y el problema del fin de semana del volumen del
tercer año sobre la fórmula de Stirling, donde la contabilidad de
suma frente a \omterm{thm:b1:integration:def}{integral}, llevada un orden más lejos, produce el propio
$\sqrt{2\pi}$.
\end{solution}
